Un ordinateur est capable de faire des {\em additions} $(15 + 3) = 18$, des {\em soustractions} 
$(15 - 3) = 12$, 
des {\em multiplications} $(15 \ast 3) = 45$ et des {\em divisions} d'entiers $(15 \;/\; 3) = 5$. 
Si la division ne tombe pas juste,
le résultat est le quotient entier. Le reste est donné
par l'opération {\em modulo} notée {\mcl mod}. Par exemple 
$(17 \;/\; 3) = 5$ et $(17 \;${\mcl mod}$\; 3) = 2$ car la division du dividende $17$
par le diviseur $3$ donne le quotient $5$ et le reste $2$. En mathématiques on utilise une croix
`$\times$' pour noter la multiplication tandis qu'en \ocaml\ on utilise l'astérisque $\ast$.
On évite ainsi de confondre le signe multiplié `$\times$' avec la lettre {\idt x} minuscule.
De même on utilise la barre oblique `$/$' pour la division entière plutôt que le signe
`$\div$' qui ne figure pas sur les claviers d'ordinateurs.
Nous verrons plus tard qu'\ocaml\ utilise d'autres opérateurs pour calculer avec des 
nombres à virgule.

Pour faire un calcul comportant plusieurs opérations, on utilise des pa\-ren\-thè\-ses. Par exemple
$((2 \ast 3) - 1)$ est égal à 5. Le calcul commence par l'évaluation de l'expression 
entre parenthèses la plus intérieure. 
Par exemple $((((5 + 1) \ast 3) - 4) \;/\; 2) = (((6 \ast 3) - 4) \;/\; 2) = 
((18 - 4) \;/\; 2) = (14 \;/\; 2) = 7$.  

\index{Expression entière}Dans ces {\em expressions entières}
nous pouvons également utiliser la valeur des compteurs de boucles ``for''. Par exemple 
lors des itérations successives dans le corps de la boucle
``{\idt  {\mcl for}\ i\ =\ 1\ {\mcl to}\ 5\ {\mcl do}\ \ldots\ {\mcl done}}'',
le compteur {\idt i} vaut $1$, puis $2$, puis $3$, puis $4$ et enfin $5$. 
Dans ce corps de boucle l'expression $((2 \ast ${\idt  i}$) - 1)$ vaut donc successivement 
$((2 \ast 1) - 1) = 1$, $((2 \ast 2) - 1) = 3$, $((2 \ast 3) - 1) = 5$, $((2 \ast 4) - 1) = 7$ et enfin
$((2 \ast 5) - 1) = 9$.

\begin{exemple}[Pyramide]
Pour dessiner la pyramide ci-après, nous pouvons commencer par le sommet et dessiner de
haut en bas des rectangles successifs de longueurs 1, 3, 5, 7, 9. La longueur du
côté du {\idt i}$^{\rm\grave{e}me}$ rectangle à partir du sommet est 
donc $((2\;\ast\;${\idt  i}$) - 1)$. Dans le programme, cette expression est écrite
après le mot {\mcl to} pour indiquer le nombre d'itérations de la boucle qui dessine un côté
du rectangle.

\pagebreak[3]
\vspace*{22mm}%
\includePostScript{Pyramide}{0.95}{0.95}%
\vspace*{-22mm}%
\input ocaml/Pyramide.tex
\end{exemple}

%\pagebreak
\noindent \rule{137mm}{0.1mm} %---------------------------------------------
\begin{exercice}{ Compléter le programme {\idt PiquetsEnDentsDeScie}, en remplaçant les
points de suspension `{\idt ...}' par une expression entière, pour
ob\-te\-nir l'alignement de piquets représenté ci-après~:
\clearpage

\raggedright{
\vspace*{20mm}%
\includePostScript{PiquetsEnDentsDeScie}{0.95}{0.95}\ \\%
}


\input ocaml/PiquetsEnDentsDeScie.tex
\noindent $\Box$}\end{exercice}
\begin{exercice}{ Écrire des programmes \ocaml\ de commande du \robot\ pour réaliser les
dessins géométriques de la figure  \ref{fig:spcf} (page  \pageref{fig:spcf}). $\Box$}\end{exercice}
\begin{exercice}{{\bf (Difficile)} Compléter le programme {\idt VaguesDePiquets} pour
des\-si\-ner l'alignement de piquets représenté ci-dessous~:

\pagebreak[3]
\raggedright{
\vspace{0.5cm}

\includevspacePostScript[2.0cm]{VaguesDePiquets}{0.95}{0.95}
}

\input ocaml/VaguesDePiquets.tex
\noindent $\Box$}\end{exercice}
\begin{figure}[hb]%[htbp]
 \includevspacePostScript[9.3cm]{SpiralesEtQuad}{0.95}{0.95}
\caption{Grecque, spirales, pavage triangulaire et carré à franges}
 \label{fig:spcf}
\end{figure}

\clearpage

\begin{exercice}{{\bf (Difficile)}
\begin{enumerate}
\item Écrire un programme \ocaml\ pour dessiner le pavage de Nicomaque\@%
\footnote{\ Nicomaque de Gérase, mathématicien grecque, env. I$^{\rm er}$ siècle après J.-C.}
représenté à la figure
\ref{fig:Nicomaque}  page \pageref{fig:Nicomaque} en utilisant des boucles
``for'' imbriquées dépendantes. Décrire les invariants de
boucles par des commentaires dans le programme. Le pavage de Nicomaque est obtenu par 10 
couronnes constituées de carrés dont les longueurs des côtés
sont respectivement 1, 2, 3, 4, \ldots, 10.

\item Calculer la surface du carré de deux façons différentes
suggérées par la figure \ref{fig:Nicomaque2}  
page \pageref{fig:Nicomaque2}~:
 \begin{enumerate}
  \item en calculant la longueur du demi-côté du carré obtenue en additionnant
        celles des carrés du pavage,
  \item en calculant la surface du quart du grand carré délimité par
   les diagonales de ce grand carré (voir cette partie en vert à la
   figure \ref{fig:Nicomaque2}, page \pageref{fig:Nicomaque2}), la surface de ce quart de grand carré
   étant obtenue par addition des surfaces des petits carrés le constituant.
 \end{enumerate}
\item En déduire l'égalité~:
\[ 1^{3} + 2^{3} + 3^{3} + 4^{3} + \ldots + 10^{3} = (1 + 2 + 3 + 4 + \ldots + 10)^{2}\]
\end{enumerate}
\noindent $\Box$}\end{exercice}
\begin{figure}[htbp]
 \includevspacePostScript[11.7cm]{Nicomaque}{0.95}{0.95}
 \caption{Pavage de Nicomaque (Échelle $\frac{1}{4})$}
 \label{fig:Nicomaque}
\end{figure}

\begin{corrige} {L'expression entière manquante dans le programme {\idt Piquets\-En\-Dents\-De\-Scie} est~:
\begin{center}
$(1 + ((${\idt i}$ - 1)\;${\mcl mod}$\;5))$
\end{center}
qui prend successivement les valeurs $1$, $2$, $3$, $4$, $5$, $1$, $2$, $3$, \dots quand {\idt i} prend les
valeurs $1$, $2$, $3$, $4$, $5$, $6$, $7$, $8$, \ldots
\noindent $\Box$}\end{corrige}
\begin{corrige} {Nous donnons un programme, les autres se trouvant dans le dossier
{\idt livre/\allowbreak 08\_expressions\_entieres}~:
\input ocaml/SpiraleTrapezoidale.tex
\noindent $\Box$}\end{corrige}
\begin{corrige} {L'expression entière manquante dans le programme {\idt VaguesDePiquets} est la suivante~:
\begin{center}
$(((6-${\idt j}$)\ast ((${\idt i}$+1)\;${\mcl mod}$\;2))+(${\idt j}$\ast(${\idt i}$\;${\mcl mod}$\;2)))$
\end{center}
Quand {\idt i} est impair, $((${\idt i}$+1)\;${\mcl mod}$\;2) = 0$ et $(${\idt i}$\;${\mcl mod}$\;2) = 1$
donc l'expression vaut {\idt j}. Quand {\idt i} est pair, $((${\idt i}$+1)\;${\mcl mod}$\;2) = 1$ et 
$(${\idt i}$\;${\mcl mod}$\;2) = 0$ donc l'expression vaut $(6-${\idt j}$)$.
$\Box$}\end{corrige}
\begin{corrige}{\hfill
\input ocaml/PavageDeNicomaque.tex
\begin{figure}[htbp]
 \includevspacePostScript[12.1cm]{Nicomaque2}{0.95}{0.95}
 \caption{Côté et surface du carré de Nicomaque (Échelle $\frac{1}{4}$)}
 \label{fig:Nicomaque2}
\end{figure}
En partageant un côté du grand carré en deux, la longueur de ce
demi-côté est la somme $(1 + 2 + \ldots + 10)$ des longueurs des carrés 
de tailles 1, 2, 3, \ldots, 10 du pavage. Ceci apparaît clairement sur la
figure \ref{fig:Nicomaque2}  page \pageref{fig:Nicomaque2} grâce à la
réglette que nous avons placée en dessous à droite du grand carré.
La surface du grand carré est donc le carré de la longueur du côté soit~:
\[(2 \times (1 + 2 + \ldots + 10))^{2} = 4 \times (1 + 2 + \ldots + 10)^{2}\]

Si l'on partage le grand carré selon les diagonales, on obtient un
triangle comme on le voit en vert à
la figure \ref{fig:Nicomaque2} (page \pageref{fig:Nicomaque2}). Ce triangle
contient 1 carré de longueur de côté 1, 2 carrés de longueur de 
côté 2, \ldots, 10 carrés de longueur de côté 10. Sa surface est
donc $1 \times 1^{2} + 2 \times 2^{2} + \ldots + 10  \times  10^{2}$ soit
$1^{3} + 2^{3} + \ldots + 10^{3}$. En faisant pivoter 4 fois ce triangle, on
recouvre entièrement le grand carré. La surface du grand carré est 
donc~:
\[4 \times (1^{3} + 2^{3} + \ldots + 10^{3})\]

Le résultat est le même quelle que soit la façon de calculer la surface 
du carré de Nicomaque. Par conséquent~:
\[4 \times (1^{3} + 2^{3} + \ldots + 10^{3}) = 4 \times (1 + 2 + \ldots + 10)^{2}\]
soit~:
\[ 1^{3} + 2^{3} + \ldots + 10^{3} = (1 + 2 + \ldots + 10)^{2}\]

Une autre démonstration de cette identité a été proposée par Al-Karagi\/\@\footnote{\ Al-Karagi,
fin X$^{\rm\grave{e}me}$ -- début XI$^{\rm\grave{e}me}$ siècle, mathématicien persan.} et repose sur le pavage du carré en 
équerres comme illustré ci-dessous~:

\begin{figure}[htbp]
 \includevspacePostScript[6.3cm]{AlKaragi}{0.95}{0.95}
 \caption{Pavage de Al-Karagi}
 %\label{fig:AlKaragi}
\end{figure}

\noindent L'aire de l'équerre extérieure de petit côté 10 vaut~:
\begin{center}
$\begin{array}{ll}
(2 \times 10 \times (1 + 2 + \ldots + 10))  - 10^{2} & = 2 \times 10 \times \frac{10 \times (10 +1)}{2} - 10^{2}\\
& = -10^{3} + 10^{2} - 10^{2} = 10^{3}
\end{array}$
\end{center}
L'aire de l'équerre de petit côté 9 est $9^{3}$, celle de l'équerre de petit côté 8 est $8^{3}$,  \ldots, 
l'aire du petit carré est $1 = 1^{3}$. 
Par conséquent, l'aire du grand carré qui
est la somme des aires du petit carré et des équerres est $1^{3} + 2^{3} + \ldots + 10^{3}$.
D'autre part, la longueur du côté du grand carré est $(1 + 2 + \ldots + 10)$ donc sa surface est
$(1 + 2 + \ldots + 10)^{2}$. Comme l'aire du grand carré est la même quelque soit la façon 
de la calculer, on obtient~:
\[ 1^{3} + 2^{3} + \ldots + 10^{3} = (1 + 2 + \ldots + 10)^{2}\]
Pour justifier que $(1 + 2 + \ldots + 10) = \frac{10 \times (10 +1)}{2}$, on utilise un rectangle de
largeur 10 et de longueur  (10 +1) dont la moitié est pavée comme suit~:

\begin{figure}[htbp]
 \includevspacePostScript[4.3cm]{S1p2p3Etc}{0.95}{0.95}
 \caption{Cacul de $1 + 2 + 3 + \ldots + 10$}
 %\label{fig:1+2+3etc}
\end{figure}

Une autre justification géométrique de l'égalité 
$1^{3} + 2^{3} + \ldots + 10^{3} = (1 + 2 + \ldots + 10)^{2}$ proposée par S. Go\-lomb
\cite{golomb1} est illustrée ci-dessous~: 
%à la figure \ref{fig:golomb} (page \pageref{fig:golomb}).

\begin{figure}[htbp]
 \includevspacePostScript[6.0cm]{Golomb}{0.95}{0.95}
 \caption{Pavage de Golomb}
 \label{fig:golomb}
\end{figure}

\noindent La longueur du côté du grand carré est bien $(1 + 2 + \ldots + 10)$ donc sa surface est
$(1 + 2 + \ldots + 10)^{2}$. Cette surface est égale à la somme des surfaces des carrés
constituant le pavage soit 1 carré de côté 1, 2 carrés de côté 2, \ldots, 10
carrés de côté 10. On 
remarque que pour les carrés de longueur de côté paire montrés en jaune sur la figure, 
il y a un 
recouvrement de deux petits carrés figuré en rouge. Mais chacun de ces petits carrés 
est adjacent à une région non couverte de même taille montrée en blanc. On peut donc
enlever un des deux petits carrés qui se recouvrent pour boucher le trou et obtenir un pavage parfait.
 La surface est donc égale à $1^{3} + 2^{3} + \ldots + 10^{3}$. 
\noindent $\Box$}\end{corrige}

%%%%%%%%%%%%%%%%%%%%%%%%%%%%%%%%%%%%%%%%%%%%%%%%%%%%%%%%%%%%%%%%%%%%%%%%%%%
\noindent \framebox[\textwidth]{
\begin{minipage}{\myparboxwidth}
\setlength{\parindent}{2em}
\small \sf
\vspace{\abovedisplayskip}
\section*{Addenda}
\indent \index{Expression entière}Pour simplifier l'apprentissage de la syntaxe d'\ocaml\ nous écrivons les expressions 
entières avec des parenthèses.  Les expressions entières {\em unaires} (avec
un seul {\em opérande} $E$) sont de la forme $(-E)$ où $-$ est l'{\em opérateur} de signe moins.
Les expressions entières {\em binaires} (avec un {\em opérande} gauche $E_1$ et un opérande 
droit  $E_2$) sont  de la forme $(E_1\,$ {\mcl op} $\,E_2)$ où {\mcl op} est l'un des 
{\em opérateurs} d'addition $+$, de soustraction $-$, de multiplication $\ast$, de division
entière $/$ ou de reste de la division entière {\mcl mod} (modulo). 
Les opérandes $E$, $E_1$ et $E_2$ sont une cons\-tan\-te entière (positive comme $1$, $17345$
(sans blancs), \ldots
ou négative comme $-10$, $-5247$, \ldots), une variable (pour l'instant un compteur de 
boucle ``for'' comme {\idt i}) ou une
{\em sous-expression} entière de la forme que nous venons de décrire. 

Le calcul de la valeur d'une expression s'appelle l'{\em évaluation} de l'expression. Les règles
d'évaluation des expressions complètement parenthésées sont très simples. La valeur d'une
constante est l'entier qui est dénoté par cette constante. La valeur d'une variable {\idt i} utilisée
comme compteur dans le corps $C$ d'une boucle {\mcl for\ \idt i\ =\ 1\ \mcl to\ \idt 
$n$\mcl\ do\ $C$\ done} (où $n \geq 1$ est une constante entière) est successivement 1 puis 2
puis \dots puis $n$ lors des itérations successives dans la boucle. La valeur d'une expression $(-E)$
s'obtient en calculant la valeur $v$ de l'expression $E$ puis en lui appliquant le signe moins. La
valeur d'une expression $(E_1\,$ {\mcl op} $\,E_2)$ s'obtient en calculant d'abord la valeur 
$v_1$ de l'expression $E_1$ puis en calculant la valeur $v_2$ de l'expression $E_2$ puis en
calculant la valeur $v_1\,$ {\mcl op} $\,v_2$ du résultat. Les mêmes règles d'évaluation
s'appliquent aux sous-expressions $E$, $E_1$ et $E_2$.

\indent Pour éviter les lourdeurs d'écriture, on peut suivre en \ocaml\ la pratique courante en
mathématiques qui consiste à utiliser des {\em règles de priorité} (on dit aussi {\em règles
de précédence}) permettant d'éliminer
certaines parenthèses redondantes.  
\vspace{\abovedisplayskip}
\end{minipage}
}
%%%%%%%%%%%%%%%%%%%%%%%%%%%%%%%%%%%%%%%%%%%%%%%%%%%%%%%%%%%%%%%%%%%%%%%%%%%%
%%%%%%%%%%%%%%%%%%%%%%%%%%%%%%%%%%%%%%%%%%%%%%%%%%%%%%%%%%%%%%%%%%%%%%%%%%%
\noindent \framebox[\textwidth]{
\begin{minipage}{\myparboxwidth}
\setlength{\parindent}{2em}
\small \sf
\vspace{\abovedisplayskip}
Les règles d'évaluation des expressions entières sont
définies et illustrées ci-après~:
\begin{list}{$\scriptscriptstyle\bullet$}{
}
\item En \ocaml\ comme en mathématiques, l'expression
entière $10 - 10 - 10$ vaut $((10 - 10) - 10)$ soit $-10$.  De même $ 2 \ast 7 \,/
\, 3\,${\mcl mod}$\, 2$ vaut $(((2 \ast 7) \,/\, 3)\,${\mcl mod}$\, 2)$ soit $0$. 
Dans les deux cas la règle est que des opérateurs de même {\em priorité} (à savoir $+$ et 
$-$ d'une part$\:$; $\ast$, $/$ et {\mcl mod} d'autre part) s'évaluent de gauche à droite. 
\item L'expression $8 + - 2 \ast 3$ vaut $(8 + ((-2) \ast 3))$ soit $2$.  De même 
\hbox{$1 - 5\,/\,2$}
vaut $(1 - (5\,/\,2))$ soit $-1$. La règle est qu'en présence d'opérateurs de
priorités différentes, on commence par évaluer les opérateurs de plus forte priorité. 
En \ocaml\ les priorités sont les suivantes~:
\end{list}

\vspace{-6pt}

\begin{center}
\hspace{1.0cm}\begin{tabular}{|l|l|l|} \hline
\multicolumn{3}{|l|}{\em Priorité des opérateurs entiers} \\ \hline \hline
$-$ (unaire)                                & opérateurs unaires & priorité forte \\ \hline
$\ast$\ $/$\ {\mcl mod} & opérateurs multiplicatifs & priorité moyenne \\ \hline
$+$, $-$ (binaires) & opérateurs additifs & priorité faible \\ \hline
\end{tabular}
\end{center}

\vspace{-6pt}

\begin{list}{$\scriptscriptstyle\bullet$}{
}
\item L'expression $(8 + - 2) \ast 3$ vaut $((8 + (- 2)) \ast 3)$ soit $18$.  De même 
\hbox{$(1 - 5)\,/\,2$} vaut $((1 - 5)\,/\,2)$ soit $-2$. Dans les deux cas on 
commence par l'évaluation des sous-expressions entre parenthèses.
\end{list}

\smallskip

\noindent De manière équivalente, l'ordre d'évaluation des expressions peut être défini
en indiquant à quels opérateurs sont liés les opérandes d'une expression~:
\begin{list}{$\scriptscriptstyle\bullet$}{
}
\item Un opérande figurant entre deux opérateurs de priorités différentes est lié à l'opérateur
de plus forte priorité (donc $2 + 3 \ast 5 = (2 + (3 \ast 5))$).
\item Un opérande figurant entre deux opérateurs de même priorité est lié à l'opérateur
de gauche (donc $2 - 3 - 5 = ((2 - 3) - 5)$).
\item Les expressions entre parenthèses sont évaluées avant d'être utilisées comme opérandes.
\end{list}

\smallskip

Les règles d'évaluation des expressions changent d'un langage de programmation à l'autre.
Par exemple en {\sc apl} \cite{iverson1}, il n'y a pas de priorités d'opérateurs (car il y en a
trop dans le langage pour qu'on puisse espérer s'en souvenir) et les expressions sont associatives 
à droite de sorte que $10 - 10 - 10 = (10 - (10 - 10)) = (10 - 0) = 10$ . Cette règle contraire
à la tradition mathématique a l'avantage que l'analyse grammaticale des expressions est grandement
simplifiée pour l'ordinateur. 

Les langages de programmation plus classiques qu'{\sc apl} sui\-vent la convention
mathématique d'évaluation de gauche à droite avec priorités.


\vspace{\abovedisplayskip}
\end{minipage}
}
%%%%%%%%%%%%%%%%%%%%%%%%%%%%%%%%%%%%%%%%%%%%%%%%%%%%%%%%%%%%%%%%%%%%%%%%%%%%
%%%%%%%%%%%%%%%%%%%%%%%%%%%%%%%%%%%%%%%%%%%%%%%%%%%%%%%%%%%%%%%%%%%%%%%%%%%
\noindent \framebox[\textwidth]{
\begin{minipage}{\myparboxwidth}
\setlength{\parindent}{2em}
\small \sf
\vspace{\abovedisplayskip}
Cependant certains langages possèdent beaucoup d'opérateurs.
Par exemple le langage {\sc C} \cite{kernighan1} comprend 43 opérateurs différents répartis en
15 niveaux de prio\-ri\-té.
Les priorités des opérateurs ne sont donc pas faciles à retenir. 

De plus les règles changent d'un langage à 
l'autre ce qui rend la vie difficile aux informaticiens
pratiquant journellement plusieurs langages. Pour ne pas se tromper ils utilisent des parenthèses
redondantes,
ce que nous conseillons aux débutants de faire.

Signalons enfin qu'en \ocaml\ la valeur des expressions entières doit être comprise entre deux
valeurs minimale et maximale, appelées {\idt min\_int} et {\idt max\_int}, qui sont
$-4611686018427387904$ et $4611686018427387903$ sur les ordinateurs actuels (elles étaient
$-32768$ et $32767$ sur les premiers ordinateurs personnels). Au-delà, le calcul donne un
résultat faux sans prévenir~: par exemple $(${\idt max\_int}$ + 1)$ vaut {\idt min\_int}$\:$!

\vspace{\abovedisplayskip}
\end{minipage}
}
%%%%%%%%%%%%%%%%%%%%%%%%%%%%%%%%%%%%%%%%%%%%%%%%%%%%%%%%%%%%%%%%%%%%%%%%%%%%
